
\documentclass[a4paper,10pt]{article}

\usepackage[utf8]{inputenc}
\usepackage[T1]{fontenc}
\usepackage{geometry}
\usepackage{booktabs}
\usepackage{array}
\usepackage{tabularx}
\usepackage{longtable}
\usepackage{listings}
\usepackage{xcolor}
\usepackage[hidelinks]{hyperref}
\usepackage{amsmath}

% Türkçe babel paketinin "=" karakterini listings seçeneklerinde etkilemesini engeller

\geometry{
    a4paper,
    top=1.8cm,
    bottom=1.8cm,
    left=1.8cm,
    right=1.8cm
}

\setlength{\parindent}{0pt}
\setlength{\parskip}{5pt}
\renewcommand{\arraystretch}{1.18}


\lstdefinestyle{cstyle}{
    language={[ANSI]C},
    basicstyle=\ttfamily\small,
    keywordstyle=\bfseries,
    commentstyle=\itshape,
    numbers=left,
    numberstyle=\tiny,
    numbersep=7pt,
    frame=single,
    breaklines=true,
    showstringspaces=false,
    tabsize=4
}

\newcommand{\hex}[1]{\texttt{0x#1}}
\newcommand{\code}[1]{\texttt{#1}}

\title{\textbf{DEOS Core Protocol Architecture}\\
\large Sürüm 1.1 -- Medya Bağımsız Çekirdek Protokol}
\author{DEOS -- Dokuz Eylül Otonom Sistemler}
\date{12 Ekim 2026}

\begin{document}
\maketitle

\begin{center}
\begin{tabular}{ll}
\textbf{Doküman tipi:} & Core Protocol Specification \& Architecture ICD \\
\textbf{Protokol sürümü:} & v1.1 (\hex{11}) \\
\textbf{Taşıma altyapısı:} & Medya Bağımsız (CAN-FD, Ethernet, LoRa, ROS 2) \\
\textbf{Durum:} & ICD Sürüm 1.1 \\
\end{tabular}
\end{center}


\section{Amaç ve Kapsam}

Bu doküman, Dokuz Eylül Otonom Sistemler (DEOS) dağıtık araç sisteminde kullanılan \textbf{medya bağımsız çekirdek protokolü} tanımlar.

DEOS protokolü yalnızca belirli bir donanıma (CAN-FD vb.) bağlı bir hat protokolü değildir; araç kontrol bilgisayarları (STM32), telemetri birimleri (LoRa), gömülü Linux / Raspberry Pi sistemleri (Ethernet TCP/UDP) ve otonom sürüş yazılımları (ROS 2) arasında ortak çalışan üst düzey bir uygulama ve taşıma mimarisidir.

Bu dokümanda:
\begin{itemize}
    \item Mantıksal Düğüm Adresleme (Node Address Registry),
    \item Evrensel Mesaj Öncelik Mimarisi (Priority Registry),
    \item Mesaj Davranış Sınıfları (Message Class Registry),
    \item Fonksiyonel Servisler (Service Registry ve Node--Service İlişkisi),
    \item Medya Bağımsız Evrensel Taşıma Çerçevesi (Universal DEOS Transport Frame),
    \item İstek--Cevap (Request--Response) ve Konfigürasyon semantiği
\end{itemize}
tanımlanmaktadır.

Fiziksel taşıma ortamlarına özel donanımsal kodlama ve çerçeveleme ayrıntıları (CAN-FD için 29-bit CAN ID ve DLC yuvarlama, LoRa için UART senkronizasyonu ve CRC, Ethernet için TCP akış çerçeveleme) ilgili \textbf{Taşıma Adaptörü (Transport Adapter)} dokümanlarında tanımlanmıştır.

Bir Node ID fiziksel bir mikrodenetleyiciye birebir bağlı olmak zorunda değildir. Aynı fiziksel cihaz üzerinde birden fazla mantıksal DEOS Node ID bulunabilir; her Node ID sistemde bağımsız bir mantıksal protokol kimliğidir.

\section{Namespace Hiyerarşisi}

DEOS protokolünde bir mesajın kimliği tek bir ``mesaj numarası'' ile değil, katmanlı bir namespace bileşimiyle belirlenir:

\begin{verbatim}
DEOS_PROTOCOL
|
+-- NODE
|   +-- SOURCE_NODE
|   +-- DESTINATION_NODE
|
+-- PRIORITY
|
+-- MESSAGE_CLASS
|
+-- SERVICE
|   +-- COMMAND
|       +-- PAYLOAD
|           +-- PARAMETER_ID   (gerektiğinde)
|           +-- VALUE          (gerektiğinde)
|
+-- PROTOCOL_METADATA
    +-- VERSION
    +-- SEQUENCE
    +-- PAYLOAD_LENGTH
\end{verbatim}

Buna göre bir mesajın semantik kimliği genel olarak şu bileşenlerden oluşur:

\begin{center}
\code{Priority + Message Class + Service + Destination + Source + Command}
\end{center}

Parametre yönetimi yapılan mesajlarda buna ek olarak \code{Parameter ID} kullanılır.

\section{Mesaj Öncelik Mimarisi (Priority Registry)}

DEOS çekirdeğinde (\code{deos\_core}) ve tüm taşıma çerçevelerinde Öncelik (\code{Priority}) alanı \textbf{1 Byte (\code{uint8\_t})} olarak tanımlanır ve taşınır. Protokol standardı olarak bu 1 baytlık alan \hex{00} ile \hex{07} (0 -- 7) arasındaki değerleri alır.

Öncelik değeri; gerek donanımsal bus arbitration (CAN-FD'ye aktarılırken 29-bit ID'nin 3 bitlik alanına haritalanır) gerekse yazılımsal kuyruk (Message Queue, LoRa TX, Router) seviyesinde mesajın öncelik sırasını belirler.

Sayısal olarak daha küçük olan değer her zaman daha yüksek önceliği temsil eder:

\begin{center}
\begin{tabularx}{\textwidth}{c l X}
\toprule
Değer & Symbol & Kullanım \\
\midrule
0 & \code{DEOS\_PRIO\_EMERGENCY} & Acil durdurma ve en kritik safety mesajları (E-Stop) \\
1 & \code{DEOS\_PRIO\_CRITICAL\_CONTROL} & Kritik gerçek zamanlı kontrol \\
2 & \code{DEOS\_PRIO\_CONTROL} & Normal sürüş kontrol mesajları \\
3 & \code{DEOS\_PRIO\_STATUS} & Durum ve geri besleme mesajları \\
4 & \code{DEOS\_PRIO\_CONFIG} & Konfigürasyon ve kalibrasyon \\
5 & \code{DEOS\_PRIO\_NETWORK} & Network management / heartbeat / ping-pong \\
6 & \code{DEOS\_PRIO\_LOW} & Düşük öncelikli servis trafiği \\
7 & \code{DEOS\_PRIO\_BACKGROUND} & Arka plan / geliştirme trafiği \\
\bottomrule
\end{tabularx}
\end{center}

\section{Node Address Registry}

Her fiziksel veya sanal DEOS haberleşme düğümüne sistem genelinde benzersiz bir 8-bit
Logical Node ID atanır. Node ID, düğümün fiziksel donanım adresi veya transport adresi
değildir; protokol içerisindeki mantıksal kimliğidir.

\subsection{Adres Aralıkları}

\begin{center}
\begin{tabularx}{\textwidth}{l X}
\toprule
Adres Aralığı & Kullanım \\
\midrule
\hex{00} & Invalid / Unassigned \\
\hex{01}--\hex{0F} & Merkezi sistem ve virtual node'lar \\
\hex{10}--\hex{1F} & Aktüatör / motor ECU düğümleri \\
\hex{20}--\hex{2F} & Güç sistemi / BMS düğümleri \\
\hex{30}--\hex{3F} & Gateway / protocol bridge düğümleri \\
\hex{40}--\hex{5F} & Sensör düğümleri için ayrılmış alan \\
\hex{60}--\hex{DF} & Gelecek genişleme \\
\hex{E0}--\hex{EF} & Geliştirme / test düğümleri \\
\hex{F0}--\hex{FD} & Diagnostik ve servis araçları \\
\hex{FE} & CAN Broadcast (Sadece CAN arayüzü) \\
\hex{FF} & Global Broadcast (Tüm arayüzler) \\
\bottomrule
\end{tabularx}
\end{center}

\subsection{Virtual Node Atamaları}

Virtual node, bağımsız bir fiziksel mikrodenetleyici veya bilgisayar olmak zorunda olmayan,
ancak DEOS protokolünde kendi Source ve Destination Node ID'si ile adreslenebilen mantıksal
haberleşme ucudur.

\begin{center}
\begin{tabularx}{\textwidth}{c l l X}
\toprule
Node ID & Symbol & Çalıştığı Ortam & Görev \\
\midrule
\hex{01} & \code{DEOS\_NODE\_TEXTUAL} & Raspberry Pi / Textual & Operatör arayüzü; config, command, request/response, status ve log sink \\
\hex{02} & \code{DEOS\_NODE\_GROUND\_CONTROL} & Uzak yer kontrol uygulaması & Yer kontrol haberleşme ucu \\
\hex{04} & \code{DEOS\_NODE\_MICRO\_ROS} & Raspberry Pi / micro-ROS & ROS arayüzü, telemetry ve control entegrasyonu \\
\bottomrule
\end{tabularx}
\end{center}

Textual Node yalnızca konfigürasyon amacıyla kullanılan bir node değildir. Protokolün izin
verdiği servislerde komut gönderebilir, cevap alabilir, durum bilgilerini işleyebilir ve
log mesajlarının doğrudan hedefi olabilir. Benzer şekilde micro-ROS ve Ground Control
Node'ları da ihtiyaç duydukları servisleri kaynak veya hedef olarak kullanabilir.

\subsection{Physical Node Atamaları}

\begin{center}
\begin{tabularx}{\textwidth}{c l l X}
\toprule
Node ID & Symbol & Fiziksel Cihaz & Görev \\
\midrule
\hex{03} & \code{DEOS\_NODE\_MAIN\_STM32} & STM32H563ZI & Ana fiziksel kontrol düğümü \\
\hex{10} & \code{DEOS\_NODE\_TRACTION} & STM32 Traction ECU & Tahrik sistemi kontrolü \\
\hex{11} & \code{DEOS\_NODE\_STEERING} & STM32 Steering ECU & Direksiyon sistemi kontrolü \\
\hex{12} & \code{DEOS\_NODE\_BRAKE} & STM32 Brake ECU & Fren sistemi kontrolü \\
\hex{20} & \code{DEOS\_NODE\_BMS\_MAIN} & DALY BMS \#1 & Ana batarya BMS \\
\hex{21} & \code{DEOS\_NODE\_BMS\_AUX} & DALY BMS \#2 & Yardımcı batarya BMS \\
\hex{30} & \code{DEOS\_NODE\_BMS\_GATEWAY} & CAN Gateway & Classic CAN--CAN-FD çevirici \\
\bottomrule
\end{tabularx}
\end{center}

\subsection{Servis ve Özel Node Atamaları}

\begin{center}
\begin{tabularx}{\textwidth}{c l l X}
\toprule
Node ID & Symbol & Tür & Görev \\
\midrule
\hex{F0} & \code{DEOS\_NODE\_DIAG\_TOOL} & Servis aracı & Diagnostik ve servis erişimi \\
\hex{FE} & \code{DEOS\_NODE\_CAN\_BROADCAST} & Mantıksal & Sadece CAN-FD ağındaki düğümlere yayın \\
\hex{FF} & \code{DEOS\_NODE\_GLOBAL\_BROADCAST} & Mantıksal & CAN, LoRa, Eth dahil tüm ağ arayüzlerine yayın \\
\bottomrule
\end{tabularx}
\end{center}

\subsection{Node Enum Tanımı}

\begin{lstlisting}[style=cstyle]
typedef enum
{
    DEOS_NODE_INVALID        = 0x00,

    DEOS_NODE_TEXTUAL        = 0x01,
    DEOS_NODE_GROUND_CONTROL = 0x02,
    DEOS_NODE_MAIN_STM32     = 0x03,
    DEOS_NODE_MICRO_ROS      = 0x04,

    DEOS_NODE_TRACTION       = 0x10,
    DEOS_NODE_STEERING       = 0x11,
    DEOS_NODE_BRAKE          = 0x12,

    DEOS_NODE_BMS_MAIN       = 0x20,
    DEOS_NODE_BMS_AUX        = 0x21,

    DEOS_NODE_BMS_GATEWAY    = 0x30,

    DEOS_NODE_DIAG_TOOL      = 0xF0,

    DEOS_NODE_CAN_BROADCAST      = 0xFE,
    DEOS_NODE_GLOBAL_BROADCAST   = 0xFF,
    DEOS_NODE_BROADCAST          = 0xFF /* Alias */
} deos_node_id_t;
\end{lstlisting}


\newpage

\section{Message Class Registry}

Message Class alanı mesajın hangi amaçla gönderildiğini tanımlar. Message Class,
fonksiyonel servisten bağımsızdır.

\begin{center}
\begin{tabularx}{\textwidth}{c l X}
\toprule
ID & Symbol & Açıklama \\
\midrule
\hex{0} & \code{DEOS\_CLASS\_SAFETY} & Acil / safety-critical mesaj \\
\hex{1} & \code{DEOS\_CLASS\_COMMAND} & Normal gerçek zamanlı kontrol komutu \\
\hex{2} & \code{DEOS\_CLASS\_STATUS} & Durum / geri besleme / telemetri \\
\hex{3} & \code{DEOS\_CLASS\_EVENT} & Fault veya olay bildirimi \\
\hex{4} & \code{DEOS\_CLASS\_REQUEST} & Cevap beklenen okuma, sorgulama veya işlem isteği \\
\hex{5} & \code{DEOS\_CLASS\_RESPONSE} & Bir REQUEST veya cevap gerektiren işleme verilen protokol cevabı \\
\hex{6} & \code{DEOS\_CLASS\_CONFIG} & Parametre ve konfigürasyon işlemleri \\
\hex{7} & \code{DEOS\_CLASS\_NETWORK} & Ağ yönetimi / heartbeat \\
\hex{8}--\hex{F} & Reserved & Gelecek kullanım \\
\bottomrule
\end{tabularx}
\end{center}

\begin{lstlisting}[style=cstyle]
typedef enum
{
    DEOS_CLASS_SAFETY   = 0x0,
    DEOS_CLASS_COMMAND  = 0x1,
    DEOS_CLASS_STATUS   = 0x2,
    DEOS_CLASS_EVENT    = 0x3,
    DEOS_CLASS_REQUEST  = 0x4,
    DEOS_CLASS_RESPONSE = 0x5,
    DEOS_CLASS_CONFIG   = 0x6,
    DEOS_CLASS_NETWORK  = 0x7
} deos_message_class_t;
\end{lstlisting}

\subsection{Message Class ve Service Arasındaki Fark}

Message Class ile Service farklı sorulara cevap verir:

\begin{center}
\begin{tabularx}{\textwidth}{l X X}
\toprule
Alan & Cevapladığı soru & Örnek \\
\midrule
Message Class & Mesaj node içerisinde hangi tür işlem akışına aittir? & COMMAND, CONFIG, NETWORK, RESPONSE \\
Service & Mesaj hangi fonksiyonel alana aittir? & STEERING, BRAKE, BMS, SYSTEM \\
Command & Bu fonksiyonel alan içerisinde tam olarak hangi işlem yapılacaktır? & SET\_TARGET\_ANGLE, SET\_PARAMETER, PING \\
\bottomrule
\end{tabularx}
\end{center}

Aynı Service farklı Message Class'lar altında kullanılabilir. Örneğin STEERING servisi için:

\begin{verbatim}
COMMAND / STEERING / SET_TARGET_ANGLE
CONFIG  / STEERING / SET_PARAMETER
REQUEST / STEERING / GET_STATUS
\end{verbatim}

Burada Service her üç mesajda da direksiyon fonksiyonel alanını belirtirken Message Class,
mesajların node içerisinde farklı işlem yollarına ayrılmasını sağlar.

Tersi de mümkündür: aynı Message Class farklı Service'lerde kullanılabilir:

\begin{verbatim}
COMMAND / TRACTION / SET_TARGET_SPEED
COMMAND / STEERING / SET_TARGET_ANGLE
COMMAND / BRAKE    / SET_TARGET
\end{verbatim}

Bu örneklerde Message Class tüm mesajların normal kontrol komutu olduğunu, Service ise hangi
alt sisteme ait olduklarını belirtir.

\section{Service Registry}

Service alanı mesajın hangi fonksiyonel domaine ait olduğunu gösterir.

\begin{center}
\begin{tabularx}{\textwidth}{c l X}
\toprule
Service ID & Symbol & Amaç \\
\midrule
\hex{00} & \code{DEOS\_SERVICE\_SYSTEM} & Genel node state/mode/reset bilgileri \\
\hex{01} & \code{DEOS\_SERVICE\_TRACTION} & Tahrik sistemi \\
\hex{02} & \code{DEOS\_SERVICE\_STEERING} & Direksiyon sistemi \\
\hex{03} & \code{DEOS\_SERVICE\_BRAKE} & Fren sistemi \\
\hex{04} & \code{DEOS\_SERVICE\_BMS} & Batarya ve BMS verileri \\
\hex{05} & \code{DEOS\_SERVICE\_CALIBRATION} & Kalibrasyon işlemleri \\
\hex{06} & \code{DEOS\_SERVICE\_DIAGNOSTIC} & Fault / sağlık / diagnostik \\
\hex{07} & \code{DEOS\_SERVICE\_SENSORS} & Sensör ve IMU veri servisi \\
\hex{08} & \code{DEOS\_SERVICE\_LIGHTING} & Aydınlatma ve sinyalizasyon sistemi \\
\hex{09}--\hex{3F} & Reserved & Yeni servisler için ayrılmış alan \\
\bottomrule
\end{tabularx}
\end{center}

\begin{lstlisting}[style=cstyle]
typedef enum
{
    DEOS_SERVICE_SYSTEM      = 0x00,
    DEOS_SERVICE_TRACTION    = 0x01,
    DEOS_SERVICE_STEERING    = 0x02,
    DEOS_SERVICE_BRAKE       = 0x03,
    DEOS_SERVICE_BMS         = 0x04,
    DEOS_SERVICE_CALIBRATION = 0x05,
    DEOS_SERVICE_DIAGNOSTIC  = 0x06,
    DEOS_SERVICE_SENSORS     = 0x07,
    DEOS_SERVICE_LIGHTING    = 0x08
} deos_service_id_t;
\end{lstlisting}


\section{Node--Service İlişkisi}

Node ve Service aynı kavram değildir. Node, fiziksel veya sanal haberleşme düğümünü;
Service ise mesajın ait olduğu fonksiyonel alanı tanımlar. Bir node'un belirli bir servise
mesaj gönderebilmesi, o servisin kendisi üzerinde barındırıldığı anlamına gelmez.

Aşağıdaki tablo temel service handler sahipliğini gösterir. Her adreslenebilir DEOS node'u,
ayrıca bu ICD'de tanımlanan ortak SYSTEM/NETWORK PING--PONG davranışını desteklemek zorundadır.

\begin{center}
\begin{tabularx}{\textwidth}{l *{8}{>{\centering\arraybackslash}X}}
\toprule
Node & SYS & TRAC & STEER & BRAKE & BMS & CAL & DIAG & LIGHT \\
\midrule
TEXTUAL & X & -- & -- & -- & -- & -- & X & -- \\
GROUND\_CONTROL & X & -- & -- & -- & -- & -- & X & -- \\
MAIN\_STM32 & X & X & X & X & X & X & X & X \\
MICRO\_ROS & X & -- & -- & -- & -- & -- & X & -- \\
TRACTION & X & X & -- & -- & -- & X & X & -- \\
STEERING & X & -- & X & -- & -- & X & X & -- \\
BRAKE & X & -- & -- & X & -- & X & X & -- \\
BMS\_MAIN & X & -- & -- & -- & X & -- & X & -- \\
BMS\_AUX & X & -- & -- & -- & X & -- & X & -- \\
BMS\_GATEWAY & X & -- & -- & -- & X & -- & X & -- \\
\bottomrule
\end{tabularx}
\end{center}

Tablodaki işaretler, ilgili node'un o servise ait mesajları desteklediğini ifade eder;
mesaj gönderme yeteneğini sınırlamaz. Örneğin \code{DEOS\_NODE\_TEXTUAL}, hedefi
\code{DEOS\_NODE\_STEERING} olan ve \code{DEOS\_SERVICE\_STEERING} kullanan bir CONFIG
veya REQUEST mesajı gönderebilir.

Örneğin direksiyon düğümüne normal sürüş komutu gönderilirken
\code{Destination = DEOS\_NODE\_STEERING}, \code{Service = DEOS\_SERVICE\_STEERING}
kullanılır. Aynı fiziksel düğüme kalibrasyon komutu gönderilirken
\code{Destination = DEOS\_NODE\_STEERING}, \code{Service = DEOS\_SERVICE\_CALIBRATION}
kullanılır.

\section{Çoklu Taşıma Katmanı Mimarisi (Multi-Transport Architecture)}

DEOS protokolü, alt taşıma katmanından tamamen soyutlanmış (transport-agnostic) olarak çalışır. Protokol çekirdeği (Core Dispatcher / Router), mesajları medya bağımsız mantıksal formatta (\code{deos\_message\_t} / \code{Universal DEOS Transport Frame}) işler ve hedef düğümün adresine göre ilgili fiziksel taşıma adaptörüne yönlendirir.

Sistemde tanımlı fiziksel taşıma adaptörleri ve ilgili ICD dokümanları şunlardır:

\begin{enumerate}
    \item \textbf{CAN-FD Taşıma Adaptörü (\code{ICD\_CANFD\_Transport}):}
    Araç içi omurga haberleşmesidir. Çekirdek protokol alanları (Priority, Class, Service, Destination, Source) 29-bit Extended CAN Identifier alanına kodlanır. BRS (Bit Rate Switch) ve DLC yuvarlama kuralları bu adaptör tarafından yürütülür.
    
    \item \textbf{Ethernet TCP/UDP Adaptörü (\code{ICD\_Ethernet\_TCP}):}
    Raspberry Pi, telemetri bilgisayarı ve teşhis araçları ile yüksek bant genişlikli IP iletişimidir. Akış ayrıştırması (stream parsing) için 1 baytlık çerçeve uzunluğu (\code{Length}) öneki kullanılır.
    
    \item \textbf{LoRa UART Telemetri Adaptörü (\code{ICD\_LoRa\_Telemetry}):}
    Yer Kontrol İstasyonu ile araç arasındaki uzun mesafe kablosuz haberleşmedir. Gürültülü açık kanal nedeniyle senkronizasyon başlığı (\code{0xAA 0x55}), paket bayrakları (\code{Flag}) ve \code{CRC16-CCITT} koruma zarfı içerir.
    
    \item \textbf{ROS 2 UDP Köprüsü:}
    Otonom sürüş yazılım katmanı (Navigasyon, Algılama) ile araç kontrol düğümleri arasında çift yönlü durum ve hedef aktarımı sağlar.
\end{enumerate}

\section{Medya Bağımsız Standart DEOS Taşıma Çerçevesi (Universal Transport Frame)}

DEOS protokolü, CAN-FD donanımının 29-bit Arbitration ID ve bit paketleme kısıtlarından bağımsız olarak; Ethernet (TCP/UDP), LoRa (UART), ROS 2 köprüleri ve Yer Kontrol İstasyonu (Python) gibi tüm harici sistemlerle uçtan uca haberleşebilmek için standart bir \textbf{Evrensel Taşıma Çerçevesi (Universal DEOS Transport Frame)} tanımlar.

Dış ortamdan üretilen bir mesaj, CAN ID hesaplamasına veya donanımsal bit kaydırmalarına gerek kalmaksızın doğrudan bu standart dizilim ile paketlenir. Mesaj STM32 Ana Beyin (Router) katmanına ulaştığında, hedef CAN düğümüyse otomatik olarak 29-bit CAN ID ve CAN-FD verisine dönüştürülür.

\subsection{Standart Çerçeve Anatomisi (Byte Byte)}

Evrensel çerçeve, \textbf{9 Byte Sabit Başlık (Header)} ve ardından gelen $N$ Byte \textbf{Payload} alanından oluşur:

\begin{center}
\begin{tabularx}{\textwidth}{c l c c X}
\toprule
\textbf{Byte} & \textbf{Alan Adı} & \textbf{Tip} & \textbf{Boyut} & \textbf{Açıklama} \\
\midrule
\textbf{0} & \code{Destination Node} & \code{uint8\_t} & 1 Byte & Hedef düğüm kimliği (Örn: \hex{03} Main STM32, \hex{10} Traction, \hex{FF} Broadcast). \\
\textbf{1} & \code{Source Node} & \code{uint8\_t} & 1 Byte & Kaynak düğüm kimliği (Örn: \hex{02} Ground Control, \hex{06} Textual UI). \\
\textbf{2} & \code{Priority} & \code{uint8\_t} & 1 Byte & Mesaj önceliği (\hex{00} CRITICAL .. \hex{05} NETWORK). \\
\textbf{3} & \code{Version} & \code{uint8\_t} & 1 Byte & DEOS protokol sürümü (Varsayılan: \hex{11} = v1.1). \\
\textbf{4} & \code{Message Class} & \code{uint8\_t} & 1 Byte & Mesaj sınıfı (\hex{01} COMMAND, \hex{02} RESPONSE, \hex{07} NETWORK vb.). \\
\textbf{5} & \code{Service ID} & \code{uint8\_t} & 1 Byte & Fonksiyonel servis (\hex{00} SYSTEM, \hex{01} TRACTION, \hex{08} LIGHTING vb.). \\
\textbf{6} & \code{Command ID} & \code{uint8\_t} & 1 Byte & Servise ait komut (\hex{01} TURN\_ON, \hex{02} TURN\_OFF, \hex{08} PING vb.). \\
\textbf{7} & \code{Sequence} & \code{uint8\_t} & 1 Byte & Mesaj takip sayacı (\hex{00} -- \hex{FF}). \\
\textbf{8} & \code{Payload Length} & \code{uint8\_t} & 1 Byte & Arkadan gelen dinamik veri bayt sayısı ($N$, 0--60 byte). \\
\midrule
\textbf{9..} & \textbf{Payload} & \code{uint8\_t[$N$]} & $N$ Byte & Gerçek komut/telemetri verisi (Örn: Far ID, Ping ID, Hız hedefi). \\
\bottomrule
\end{tabularx}
\end{center}

\textbf{Toplam Paket Boyutu:} $9 + N$ Byte.

\subsection{CAN-FD Dönüşüm Mekanizması}

Bu çerçeve STM32 Gateway/Router düğümüne ulaştığında:
\begin{enumerate}
    \item \textbf{29-bit CAN ID:} Byte 0 (\code{Destination}), Byte 1 (\code{Source}), Byte 2 (\code{Priority}), Byte 4 (\code{Message Class}) ve Byte 5 (\code{Service}) alanları sürücü tarafından otomatik olarak 29-bit Extended CAN Identifier'a dönüştürülür.
    \item \textbf{CAN-FD Data Payload:} Byte 3 (\code{Version}), Byte 7 (\code{Sequence}), Byte 6 (\code{Command}), Byte 8 (\code{Payload Length}) ve ardından gelen \code{Payload} alanları doğrudan CAN-FD frame'inin veri alanına yazılır.
\end{enumerate}

\subsection{Fiziksel İletim Katmanlarına Göre Dış Zarf (L2 Framing)}

Evrensel 9+N baytlık standart DEOS çerçevesi, fiziksel iletim hattının özelliklerine göre dış zarflarla taşınır:
\begin{itemize}
    \item \textbf{Ethernet (TCP):} TCP zaten akış kontrolü ve sağlama toplamı sağladığından, çerçevenin başına sadece toplam paket boyutu eklenir: \code{[Length (1B)] + [Standart DEOS Frame (9+N Byte)]}.
    \item \textbf{LoRa / Seri Port (UART):} Gürültülü ve bayt kayması olabilecek hatlarda senkronizasyon ve CRC doğrulaması eklenir: \code{[0xAA 0x55] [Flag] [Length] + [Standart DEOS Frame (9+N Byte)] + [CRC16 (2B)]}.
    \item \textbf{ROS 2 / UDP:} Doğrudan \code{[Standart DEOS Frame (9+N Byte)]} datagram olarak iletilir.
\end{itemize}

\section{REQUEST--RESPONSE ve CONFIG Ayrımı}

REQUEST, RESPONSE ve CONFIG birbirinden bağımsız mesaj sınıflarıdır ve farklı amaçlara hizmet eder.

\begin{itemize}
    \item \code{REQUEST}: Bir düğümden veri, durum veya belirli bir işlem sonucu talep etmek için kullanılır.
    \item \code{RESPONSE}: REQUEST sınıfındaki bir mesaja veya cevap gerektiren başka bir işleme karşılık dönen cevabı taşır.
    \item \code{CONFIG}: Bir düğümün çalışma parametrelerini, ayarlarını veya kalıcı yapılandırmasını değiştirmeye yönelik mesajları taşır.
\end{itemize}

Örneğin bir parametreyi okumak için REQUEST, aynı parametreyi değiştirmek için CONFIG
kullanılabilir:

\begin{verbatim}
REQUEST / STEERING / GET_PARAMETER / PID_KP
CONFIG  / STEERING / SET_PARAMETER / PID_KP
\end{verbatim}

Her iki işlem için de gerektiğinde RESPONSE sınıfı ile sonuç döndürülebilir. Bu ayrım,
mesajın ``cevap bekleyip beklemediği'' ile ``konfigürasyon amacı taşıyıp taşımadığı''
bilgilerinin birbirine karıştırılmasını önler.

\section{Response Altyapısı}

RESPONSE Message Class protokolde ayrılmıştır; ancak hangi komutların zorunlu response
üreteceği bu registry dokümanında kesinleştirilmemiştir.

Özellikle aşağıdaki işlemler gelecekte response kullanılabilecek adaylardır:
\begin{itemize}
    \item Parameter GET/SET,
    \item state veya mode değişimi,
    \item kalibrasyon başlatma / kaydetme,
    \item fault reset / clear,
    \item node reset,
    \item servis tipi tek-seferlik yönetim komutları.
\end{itemize}

\subsection{Generic Response Payload}

\code{DEOS\_CLASS\_RESPONSE} mesajları varsayılan olarak standart bir veri yerleşimi kullanır. Eğer ilgili komut özel bir ek veri döndürmüyorsa, yanıt mesajının içeriği her zaman şu formatta olur:

\begin{center}
\begin{tabularx}{\textwidth}{c l c X}
\toprule
Byte & Alan & Tip & Açıklama \\
\midrule
0 & Result Code & \code{uint8\_t} & İşlem sonucu (\code{deos\_result\_t}) \\
1..N & Ek Veri (Opsiyonel) & Değişken & Komuta özgü veri \\
\bottomrule
\end{tabularx}
\end{center}

Result code \hex{00} (SUCCESS) olduğunda işlem başarılı sayılır. Bu genel yapı sayesinde, her komut için özel bir response formatı tasarlamaya gerek kalmadan hata yönetimi standartlaştırılmıştır.

Aşağıda standart response result kodları listelenmiştir:

\begin{center}
\begin{tabular}{c l}
\toprule
Kod & Anlam \\
\midrule
\hex{00} & SUCCESS \\
\hex{01} & UNSUPPORTED \\
\hex{02} & INVALID\_PARAMETER \\
\hex{03} & OUT\_OF\_RANGE \\
\hex{04} & NOT\_ALLOWED \\
\hex{05} & BUSY \\
\hex{06} & NOT\_READY \\
\hex{07} & INTERNAL\_ERROR \\
\bottomrule
\end{tabular}
\end{center}

Bu alanın kesin payload yerleşimi sonraki ICD revizyonunda mesaj bazında ayrıntılandırılacaktır.


\section{Mantıksal Node Kimliği Kuralı}

CAN hattında her Source ve Destination Node ID bağımsız bir DEOS protokol kimliğidir.
Aynı fiziksel cihaz üzerinde birden fazla mantıksal Node ID bulunabilse bile, bir mesaj
hangi Node ID adına üretiliyorsa CAN Identifier içerisindeki Source alanında o Node ID
kullanılmalıdır.

Bir node başka bir node'un kimliğiyle cevap üretmemelidir. Örneğin
\code{Destination = DEOS\_NODE\_BMS\_MAIN} olan geçerli bir PING'e verilen PONG mesajında
\code{Source = DEOS\_NODE\_BMS\_MAIN} kullanılmalıdır.

Bu kural fiziksel yerleşimden bağımsızdır; CAN hattındaki diğer katılımcılar mantıksal
Node ID'leri bağımsız haberleşme uçları olarak değerlendirir.

\section{İsimlendirme Standardı}

Tüm protokol sabitlerinde ortak prefix yapısı kullanılır:

\begin{verbatim}
DEOS_NODE_...
DEOS_PRIO_...
DEOS_CLASS_...
DEOS_SERVICE_...
DEOS_CMD_<SERVICE>_...
DEOS_PARAM_<SERVICE>_...
DEOS_STATE_...
DEOS_MODE_...
\end{verbatim}

Örnekler:

\begin{verbatim}
DEOS_NODE_TEXTUAL
DEOS_NODE_GROUND_CONTROL
DEOS_NODE_MAIN_STM32
DEOS_NODE_MICRO_ROS
DEOS_NODE_STEERING
DEOS_CLASS_CONFIG
DEOS_SERVICE_STEERING
DEOS_CMD_STEERING_SET_TARGET_ANGLE
DEOS_CMD_SET_PARAMETER
DEOS_PARAM_STEERING_PID_KP
DEOS_STATE_ACTIVE
DEOS_MODE_AUTONOMOUS
\end{verbatim}


\section{Özet Registry}

Mevcut CAN-FD protokol kararları aşağıdaki şekilde özetlenebilir:

\begin{itemize}
    \item Ana araç içi haberleşme hattı CAN-FD'dir.
    \item 29-bit Extended CAN Identifier kullanılır.
    \item CAN-FD frame'lerinde BRS kullanılır.
    \item Priority, Message Class, Service, Destination ve Source alanları CAN Identifier içerisinde taşınır.
    \item Node ID, sistem genelinde kullanılan 8-bit mantıksal protokol kimliğidir.
    \item Aynı fiziksel cihaz birden fazla mantıksal Node ID barındırabilir; CAN hattında her biri bağımsız kimliktir.
    \item DEOS Protocol v1.1 sürüm değeri \hex{11}'dir.
    \item Ortak DEOS data header'ı Version, Sequence, Command ve Length olmak üzere 4 byte'tır.
    \item Command-specific maksimum payload uzunluğu 60 byte'tır.
    \item Length, yalnızca geçerli command-specific payload byte sayısını belirtir.
    \item DLC fiziksel CAN-FD data field boyutunu; Length ise semantic DEOS payload boyutunu ifade eder.
    \item DLC yuvarlaması sonucu oluşan kullanılmayan byte'lar \hex{00} padding olarak gönderilir.
    \item Uygulama seviyesinde ek CRC veya ayrıca sync header kullanılmaz.
    \item Sequence alanı Source Node başına kullanılır; v1.1'de gap/duplicate kontrolü zorunlu değildir.
    \item REQUEST, RESPONSE ve CONFIG birbirinden ayrı Message Class değerleridir.
    \item CONFIG ayrı bir Service değildir.
    \item Parameter ID 16-bit'tir ve Service namespace'i içinde anlam kazanır.
    \item Varsayılan çok baytlı wire byte sırası Little-Endian'dır.
    \item Fiziksel değerlerde varsayılan yaklaşım ölçekli integer kullanımıdır.
    \item PING/PONG, \code{DEOS\_CLASS\_NETWORK / DEOS\_SERVICE\_SYSTEM} altında tanımlıdır.
    \item PING ve PONG command payload'u 4-byte Ping ID içerir ve Length alanı \hex{04}'tür.
    \item PING yalnız unicast olarak geçerlidir; Broadcast PING'e PONG üretilmez.
    \item Bir cevabın Source Node değeri, cevabı üreten mantıksal node kimliği olmalıdır.
\end{itemize}

\end{document}
